% La Brasa: Harmonised MODIS/VIIRS fire activity for the Iberian Peninsula
% NASA Space Apps Challenge 2026 -- "Harmonization of MODIS and VIIRS Hot Spots"
% Compile: pdflatex la_brasa.tex (twice). Self-contained, no external .bib needed.
\documentclass[11pt,a4paper]{article}

\usepackage[T1]{fontenc}
\usepackage[utf8]{inputenc}
\usepackage[english]{babel}
\usepackage{lmodern}
\usepackage{microtype}
\usepackage[margin=2.4cm]{geometry}
\usepackage{amsmath,amssymb}
\usepackage{booktabs}
\usepackage{tabularx}
\usepackage{array}
\usepackage{siunitx}
\usepackage{enumitem}
\usepackage{xcolor}
\usepackage[hidelinks]{hyperref}
\usepackage{caption}

\sisetup{group-separator={,},group-minimum-digits=4,output-decimal-marker={.}}
\captionsetup{font=small,labelfont=bf}
\setlist{nosep,leftmargin=*}
\newcolumntype{L}{>{\raggedright\arraybackslash}X}

\title{\textbf{La Brasa}: A harmonised, uncertainty-aware record of satellite
fire activity for the Iberian Peninsula, 2001--2026\\[4pt]
\large Technical report for the NASA Space Apps Challenge 2026,\\
challenge \emph{Harmonization of MODIS and VIIRS Hot Spots}}
\author{Team La Brasa}
\date{November 2026}

\begin{document}
\maketitle

\begin{abstract}
\noindent
Active-fire records from NASA FIRMS cover more than two decades but come from instruments with different
spatial resolution and detection sensitivity: MODIS (\SI{1}{km}, Terra and Aqua, from 2000) and VIIRS
(\SI{375}{m}, Suomi~NPP, NOAA-20, NOAA-21, from 2012). Simply joining both records produces a spurious
step in 2012, so detection counts cannot be used directly for trends or year-to-year comparison.
We present \emph{La Brasa}, an open-source pipeline and web application that harmonises all five
satellites onto a common VIIRS-equivalent scale for the 21 NUTS-2 regions of mainland Spain, the Balearic
Islands and Portugal. Conversion factors are estimated per region in the 2012--2021 sensor overlap and
validated by leave-one-year-out cross-validation against pre-registered quality thresholds; every
harmonised value carries an empirically calibrated 50\,\% and 90\,\% uncertainty band. On top of the
harmonised record we derive a burning-activity calendar, regional fire seasons, anomalies relative to
earlier years, fire events, a cloud-cover layer that marks where ``no detection'' may mean ``no view'',
and a probabilistic 1--10-day forecast of large-fire days. Before 2012 the harmonised daily count has a
typical error factor of $\times 2.08$, the monthly count $\times 1.46$; large-fire days are recovered with
Cohen's $\kappa = 0.67$. With weather inputs perturbed by real forecast errors, the forecast improves the
Brier score over seasonal climatology by 25\,\% on day~1 and 5\,\% on day~10.
\end{abstract}

% ---------------------------------------------------------------------------
\section{Introduction}

The Iberian Peninsula is one of Europe's most fire-affected regions. Emergency services, land managers
and researchers repeatedly ask three questions: \emph{Is this season unusual? Where is fire activity rising?
How far can we trust what the satellites show?} NASA's Fire Information for Resource Management System
(FIRMS) provides the relevant observations, but as individual hotspots from instruments that are not
directly comparable. VIIRS detects smaller and cooler fires than MODIS
\cite{schroeder2014,giglio2016}; in our study area Suomi~NPP reports $6.0\times$ as many detections as
MODIS-Aqua. A naive time series therefore shows an apparent surge in fire activity in 2012 that is a
measurement artefact.

The 2026 Space Apps challenge \emph{Harmonization of MODIS and VIIRS Hot Spots} asks teams to
``combine MODIS and VIIRS fire records into a consistent calendar for examining burning activity and
unusual conditions'' within a selected area of interest, in a web application for early-warning and
emergency responders, scientists and land managers. This report describes how La Brasa addresses this.
Its contributions are:
\begin{enumerate}
  \item a reproducible harmonisation of five satellites (Terra, Aqua, Suomi~NPP, NOAA-20, NOAA-21) with
        region-specific factors and a magnitude-dependent, cross-validated uncertainty model
        (Sections~\ref{sec:data}--\ref{sec:validation});
  \item derived products that turn hotspots into answers: calendar, fire season, anomalies, events and a
        cloud-obscuration layer (Section~\ref{sec:products});
  \item a probabilistic short-range forecast evaluated with real weather-forecast errors
        (Section~\ref{sec:forecast});
  \item an open web application (English, Spanish, Portuguese) built on a documented API
        (Section~\ref{sec:system}).
\end{enumerate}
Section~\ref{sec:compliance} maps every element of the challenge statement to the implementation.

% ---------------------------------------------------------------------------
\section{Data}\label{sec:data}

\paragraph{Fire detections.}
We use the FIRMS standard-processing archives MODIS Collection~6.1 (Terra and Aqua, from November 2000)
and VIIRS 375\,m Collection~2 (Suomi~NPP from 20~January~2012, NOAA-20 from 1~April~2018, NOAA-21 from
1~March~2023), extended with near-real-time (NRT) data up to mid-November 2026. One record is one pixel
flagged as thermally anomalous during an overpass; it is an observation, not a fire and not an area.

\paragraph{Study area.}
All satellites are counted only inside the 21 NUTS-2 regions (NUTS 2021 \cite{nuts2021}) of mainland Spain
with the Balearic Islands and of Portugal. Regions are assigned locally by point-in-polygon; each
detection additionally receives a $0.25^\circ$ cell identifier and an H3 cell at resolution~8
($\approx\SI{0.74}{km^2}$) \cite{h3}.

\paragraph{Filtering.}
Only vegetation fires are retained (\texttt{type}~=~0). Static land sources (industry, gas flares;
\texttt{type}~=~2) are removed. Because NRT data lack the \texttt{type} field, a location mask learned from
the archive is applied per sensor; it recovers 92.6\,\% of static detections and removes 1.5\,\% of true
fires. Of \num{940293} raw records, \num{712097} are retained; the remainder are static sources
(\num{133317}), offshore (\num{2450}), masked (\num{4274}) or outside the study regions (\num{88155}).

\paragraph{Auxiliary data.}
Weather for the forecast (ERA5 reanalysis \cite{hersbach2020} for history, ECMWF IFS and its ensemble for
forecasts) and effective cloud cover at the overpass times are retrieved through Meteomatics, with
Open-Meteo as the fallback source.

% ---------------------------------------------------------------------------
\section{Harmonisation method}\label{sec:method}

\subsection{Findings that shaped the method}
A preliminary analysis (calibration window February~2012 to December~2021) established:
\begin{itemize}
  \item \textbf{Detection threshold, not pixel size, drives the difference.} Aqua detects a VIIRS fire
        cluster with 50\,\% probability only above $\approx\SI{34}{MW}$ (day) or \SI{21}{MW} (night). Aqua
        sees 24\,\% of VIIRS clusters, which hold 83\,\% of the radiative power.
  \item \textbf{FRP is far more comparable than counts.} The FRP sum ratio SNPP/Aqua is $\approx 1.04$,
        whereas the count ratio is $6.0$ (Aqua) or $3.0$ (Aqua~+~Terra).
  \item \textbf{Simple beats complex.} Of 17 candidate conversions (monthly or day/night factors, Poisson
        models, gradient boosting, quantile mapping, FRP-class weights), a single factor per region gave the
        best out-of-sample agreement.
  \item \textbf{Orbital drift.} From 2022 the Terra and Aqua overpass times drift by more than 100 minutes.
        MODIS is therefore used only up to 2021 and only for days without VIIRS.
\end{itemize}

\subsection{Reference and conversion}
VIIRS is the reference whenever it operates. On a measured day $t$ with $k_t\in\{1,2,3\}$ active VIIRS
satellites, the reference count in region $r$ is the mean
\begin{equation}
  V_{r,t} \;=\; \frac{1}{k_t}\sum_{s\in\mathcal{S}_t} n_{s,r,t},
\end{equation}
where an active satellite without detections contributes~0. The mean, not the sum, is used because the
satellites observe the same fires about 50 minutes apart and cannot be de-duplicated; a sum would roughly
double the series from 2018. The mean stays at the SNPP level (median ratio 1.006) and is less noisy.

Before 2012 the series is converted from MODIS with region-specific factors estimated on the overlap
window $\mathcal{T}_c$ (2012--2021) against Suomi~NPP:
\begin{equation}
  f_r \;=\; \frac{\sum_{t\in\mathcal{T}_c} V^{\mathrm{SNPP}}_{r,t}}{\sum_{t\in\mathcal{T}_c} M_{r,t}},
  \qquad
  n^{\mathrm{harm}}_{r,t} =
  \begin{cases}
    V_{r,t} & \text{VIIRS operating},\\
    f^{\,\mathrm{A+T}}_r\,(n^{\mathrm{Aqua}}_{r,t}+n^{\mathrm{Terra}}_{r,t}) & \text{otherwise}.
  \end{cases}
\end{equation}
Regions with fewer than 300 MODIS detections in $\mathcal{T}_c$ fall back to the pooled factor. Fire
radiative power is converted from Aqua alone, which was more accurate than the Aqua~+~Terra pair
($R^2$ 0.59 vs.\ 0.52); before July~2002 only Terra is available. The pooled factors are 2.996 (count,
Aqua~+~Terra), 5.98 (count, Aqua) and 1.178 (FRP, Aqua). Sensors are never added together, so the series
is VIIRS-equivalent throughout. The canonical product is region~$\times$~day; weekly, monthly and yearly
values are derived by a tested aggregation module and are left empty for incomplete periods.

\paragraph{Point weights.}
To support freely drawn areas, each detection $i$ carries a weight $w_i$: $1/k_t$ for VIIRS points and
$f_r$ for MODIS points that enter the series. By construction
$\sum_{i\in(r,t)} w_i = n^{\mathrm{harm}}_{r,t}$ (verified to $\approx 10^{-16}$), so any spatial query
returns values consistent with the regional series.

\paragraph{Zero versus missing.}
A harmonised 0 means ``no detection''; an empty value means ``no admissible observation''. Before 2012 a 0
often only means that MODIS saw nothing, which the upper uncertainty band makes explicit.

\subsection{Uncertainty model}
The error of a converted value depends strongly on its magnitude. From leave-one-year-out
cross-validation (2012--2021) we estimate multiplicative error factors per spatial scale (cell, NUTS-2,
country, all), temporal scale (day to year) and magnitude class, e.g.\ for region~$\times$~month from
$\times 3.0$ for converted zeros to $\times 1.14$ above 500 detections. A single factor ($\times 1.46$)
would be far too optimistic for small values. Each value is published with 50\,\% and 90\,\% bands and a
reliability label (\emph{measured}, \emph{converted\_solid}, \emph{converted\_rough}).

% ---------------------------------------------------------------------------
\section{Validation}\label{sec:validation}

All quality indicators (KPIs) and their thresholds were fixed in the configuration \emph{before}
measurement and were never relaxed. Converted MODIS values are evaluated against held-out VIIRS years
(leave one year out, 2012--2021) on the 18 regions with at least \num{2000} VIIRS detections. All 42 KPIs
pass; Table~\ref{tab:kpi} lists the central ones.

\begin{table}[ht]
\centering\small
\caption{Selected out-of-sample KPIs of the harmonisation (converted MODIS vs.\ measured VIIRS).}
\label{tab:kpi}
\begin{tabularx}{\linewidth}{@{}lLr@{}}
\toprule
ID & Quantity & Result \\
\midrule
D1 & Typical daily error factor, region $\times$ day (90th percentile) & $\times 2.08$ ($\times 6.0$) \\
D2 & Overall count bias & 0.1\,\% \\
D3 & Monthly error factor; activity class agreement; $\kappa$ & $\times 1.46$; 63.7\,\%; 0.44 \\
D4 & Daily FRP, $R^2$ (log); ratio & 0.59; 0.975 \\
D5 & Large-fire day (FRP $\geq$ \SI{100}{MW}): $\kappa$, recall, precision & 0.67, 67\,\%, 70\,\% \\
D7 & Extreme days ($\geq$ \SI{500}{MW}) correctly classified & 76\,\% \\
D8 & Climatological season profile: $R^2$; error factor & 0.945; $\times 1.23$ \\
D10/11 & Point weights sum exactly to regional series & $\approx 10^{-15}$ \\
U1/U2 & Empirical coverage of 50\,\% / 90\,\% bands & 0.43--0.58 / 0.88--0.92 \\
M1 & Static-source mask: recall / true fires lost & 92.6\,\% / 1.5\,\% \\
N2 & VIIRS mean vs.\ SNPP alone, median ratio per region-year & 1.006 \\
\bottomrule
\end{tabularx}
\end{table}

Error shrinks with aggregation (day $\times 2.08$, week $\times 2.0$, month $\times 1.5$, year
$\times 1.12$). The consequence for users is explicit in the application: pre-2012 values are always shown
with bands, and daily pre-2012 values are suitable for identifying large-fire days and seasonal patterns,
not for exact counts. Five documented hypothesis tests (e.g.\ solar instead of UTC date, a separate
large-fire threshold for converted days) are reported with the decision taken.

% ---------------------------------------------------------------------------
\section{Derived products}\label{sec:products}

\paragraph{Burning-activity calendar and fire season.}
For any area the application shows a year $\times$ week/month heat map and a seasonal curve: the median
and inter-quartile band per calendar day across all years, smoothed over 15 days. Season start and end
are defined at 10\,\% and 90\,\% of the yearly total; the peak phase spans $\pm 7$ days around the peak.
The current year-to-date is compared with the same calendar days of earlier years (no extrapolation).

\paragraph{Anomalies.}
An anomaly map colours H3 cells by the difference to the median of the same time window in user-selected
comparison years, with value, median and rank on demand. Fire-day probabilities are computed only from
VIIRS years, because adding satellites raises the share of days with any detection (0.32 to 0.41--0.45)
without a real trend.

\paragraph{Fire events.}
Hotspots are grouped into events by union--find over H3 cells (neighbourhood ring~2, at most two days
between detections, at least two hotspots). Events carry duration, harmonised fire count, FRP, a dense core
and a three-day trend; merging events keep the older identifier.

\paragraph{Cloud-obscuration layer.}
A satellite cannot see a fire under clouds, so ``0 fires'' may mean ``no view''. For each $0.25^\circ$ cell
and day we store the effective cloud fraction at the day ($\approx$~13~UTC) and night ($\approx$~02~UTC)
VIIRS overpasses. A day is shown as cloudy when both overpasses exceed a configurable threshold (default
0.5). Fires detected on the previous and the following day but not on the day itself, at a cloudy location,
are drawn as ``probably hidden''. Fire counts are never altered. Between two fire days, a cloudy day remains
without detection 3.2 times as often as a clear one (74.3\,\% vs.\ 22.9\,\% at threshold 0.7); 22--41\,\%
of measured zero-fire region-days are cloudy, depending on the threshold. Clouds therefore hide fires
measurably, mainly outside summer, but explain only a minority of gaps. Gap filling from neighbouring days
was tested and rejected: it would add only 0.3--0.9\,\% of fires and is within a factor of two in only
55\,\% of test cases.

% ---------------------------------------------------------------------------
\section{Fire-activity forecast}\label{sec:forecast}

\paragraph{Targets and model.}
Per region, the daily product gives the probability of a large-fire day (harmonised FRP
$\geq$ \SI{100}{MW}) and of a fire day (based on Suomi~NPP, to avoid the satellite-count drift) for lead
times of 1--10 days. A weekly product gives tercile probabilities (below, normal, above usual) for weeks
2--4. Each lead time has an L2-regularised logistic regression with Platt calibration \cite{platt1999}.
Inputs are the smoothed seasonal climatology of the target day, recent fire history (log-FRP over 1, 3, 7
and 30 days, days since the last large fire, activity in the three nearest regions), a penalised regional
effect and 22 weather features: the Canadian Fire Weather Index system (FFMC, DMC, DC, ISI, BUI, FWI)
\cite{vanwagner1987}, the Fosberg index \cite{fosberg1978}, temperature, humidity, vapour-pressure deficit,
gusts, antecedent rain and their anomalies. Weather is taken at the eight $0.25^\circ$ cells per region
with the most VIIRS detections in the training years. Our FWI implementation reproduces the published test
table \cite{vanwagnerpickett1985} to within 0.005.

\paragraph{Honest evaluation.}
Training covers 2012--2021, testing 2022--2026; hyper-parameters were chosen on a validation split
(2020--2021), not on the test years. A forecast with perfect (reanalysis) weather would overstate skill.
We therefore perturb the test-period weather with real ECMWF IFS errors per lead time (Open-Meteo previous
runs against ERA5) and train on weather with such errors. Skill is the Brier skill score against seasonal
climatology, $\mathrm{BSS} = 1 - \mathrm{BS}_{\mathrm{model}}/\mathrm{BS}_{\mathrm{clim}}$ \cite{brier1950},
with 90\,\% intervals from bootstrapping over years. The model is served for a lead time only if the gain
is positive and the lower bootstrap bound is $\geq 0$; otherwise the application shows climatology and
says so.

\begin{table}[ht]
\centering\small
\caption{Forecast of large-fire days, test 2022--2026: BSS relative to seasonal climatology.}
\label{tab:forecast}
\begin{tabular}{@{}lrrrrr@{}}
\toprule
Lead time (days) & 1 & 3 & 5 & 7 & 10 \\
\midrule
Without weather & 0.213 & 0.078 & 0.037 & 0.020 & 0.005 \\
With weather, real forecast errors & \textbf{0.250} & \textbf{0.147} & \textbf{0.109} & \textbf{0.077} & \textbf{0.053} \\
\bottomrule
\end{tabular}
\end{table}

Calibration is good (expected calibration error $\leq 0.009$, reliability slope 1.01--1.13). On day~1 the
model without weather reaches AUC~0.92 against 0.83 for climatology and beats climatology in all 15
leave-one-year-out years. In the highest risk class, 54\,\% of region-days are large-fire days, against a
base rate of 2.9\,\%. Weather adds no skill to the weekly outlook, whose ranked probability skill score
\cite{wilks2011} is small but positive in week~2 (0.020) and indistinguishable from zero in week~4. The
probability range shown to users comes from the ECMWF ensemble spread.

% ---------------------------------------------------------------------------
\section{System}\label{sec:system}

The system consists of strictly separated stages that communicate only through files or the database:
\begin{enumerate}
  \item \emph{Ingestion} downloads FIRMS archive and NRT data through the FIRMS API.
  \item \emph{Harmonisation} (Python) writes versioned Parquet tables with a manifest, hashes and a KPI
        report; 54 automated tests and the KPI checks run with every full rebuild (about 90\,s).
  \item \emph{Database and API}: PostgreSQL and a FastAPI service provide daily and calendar series,
        hotspots, events, anomaly grids, cloud cover and forecasts for regions, countries or drawn polygons,
        plus CSV export with uncertainty.
  \item \emph{Web application} (React, MapLibre): area selection, H3 hex map with time slider, calendar,
        season curve, year comparison, sensor comparison (raw vs.\ harmonised), event list, cloud strip,
        forecast box and a guided tutorial, in English, Spanish and Portuguese.
\end{enumerate}
A daily job updates FIRMS data, harmonises and issues the forecast; a sync button in the application tops
up NRT data between runs. The code is released under the Apache License~2.0.

% ---------------------------------------------------------------------------
\section{Compliance with the challenge statement}\label{sec:compliance}

\begin{table}[ht]
\centering\small
\caption{Requirements derived from the challenge statement and their implementation.}
\label{tab:compliance}
\begin{tabularx}{\linewidth}{@{}p{3.6cm}LL@{}}
\toprule
Challenge wording & Requirement & Implementation \\
\midrule
``harmonizes these records'' & Continuous, comparable series; method transparent and tested &
  Regional factors, VIIRS-equivalent scale, 42 pre-registered KPIs (Sec.~\ref{sec:method}--\ref{sec:validation}) \\
``burning activity calendar'' & Calendar view of activity from hotspots &
  Year $\times$ week/month heat map, ring calendar, daily time slider \\
``selected area of interest'' & Region, country or drawn area &
  NUTS-2, country, Iberia, free polygon; point weights keep values consistent \\
``historical fire patterns'' & Typical season per area &
  Median and IQR per calendar day, season start, end and peak \\
``unusual conditions'' & Detect and flag deviations &
  Anomaly map vs.\ comparison years; year-to-date vs.\ history \\
``critical periods'' & Windows of typically high activity &
  Peak phase; forecast risk classes for the next 10 days \\
``wildfires and agricultural burns'' & Keep vegetation fires, remove industry &
  \texttt{type}~0 filter and static-source mask; wildfire vs.\ agricultural split not yet implemented \\
``web application'' & Running application, not a notebook &
  React/MapLibre app on FastAPI and PostgreSQL \\
``early warning \dots{} responders, scientists, land managers'' & Clear statements, uncertainty, no raw-data flood &
  Bands and reliability labels, cloud layer, three languages, tutorial, CSV export \\
\bottomrule
\end{tabularx}
\end{table}

% ---------------------------------------------------------------------------
\section{Limitations}

\begin{itemize}
  \item \textbf{Regional scope.} Factors are calibrated for Iberia and do not transfer: global FRP factors
        range from 0.9 to 3.3. Other regions need their own calibration and region boundaries.
  \item \textbf{Pre-2012 resolution.} MODIS misses about 60\,\% of region-days on which VIIRS detects fire.
        Daily pre-2012 values are reliable for large-fire days, not for small fires.
  \item \textbf{Measured-day bands} assume counting noise with a minimum factor; they are not yet narrowed
        for days with three VIIRS satellites, and single-satellite outages are only flagged.
  \item \textbf{Clouds} are visualised, not corrected; the layer covers 2012 onward.
  \item \textbf{Weather data} for the forecast and the cloud layer currently come from a time-limited
        commercial account; after it expires the system falls back to Open-Meteo, and the cloud layer stops
        updating.
  \item \textbf{Fire type.} Wildfires and agricultural burns are not separated yet; land-cover data would
        allow this.
\end{itemize}

% ---------------------------------------------------------------------------
\section{Conclusion}

La Brasa turns 26 years of heterogeneous satellite hotspots into one consistent, VIIRS-equivalent record
for Iberia and makes it usable in the browser. Its core design choices are: harmonise rather than merge,
publish uncertainty with every number, separate ``zero'' from ``unobserved'', and judge every model
against climatology with thresholds fixed in advance. The result answers the questions practitioners ask:
how today compares with the past, where activity is unusual, where satellites may have been blind, and
what the coming days are likely to bring.

\paragraph{Code and data availability.}
Source code, configuration, tests and method documentation are available at \url{https://github.com/Orbital-Fire-Watch/La-Brasa.git} under the
Apache License~2.0. All fire data come from NASA FIRMS; region boundaries from Eurostat GISCO.

% ---------------------------------------------------------------------------
\begin{thebibliography}{99}\small

\bibitem{brier1950}
G.~W. Brier.
\newblock Verification of forecasts expressed in terms of probability.
\newblock \emph{Monthly Weather Review}, 78(1):1--3, 1950.

\bibitem{nuts2021}
Eurostat.
\newblock \emph{NUTS -- Nomenclature of Territorial Units for Statistics, version 2021}.
\newblock GISCO, European Commission, 2021.

\bibitem{fosberg1978}
M.~A. Fosberg.
\newblock Weather in wildland fire management: the fire weather index.
\newblock In \emph{Proceedings of the Conference on Sierra Nevada Meteorology}, pp.~1--4. American
Meteorological Society, 1978.

\bibitem{giglio2016}
L.~Giglio, W.~Schroeder, and C.~O. Justice.
\newblock The collection 6 MODIS active fire detection algorithm and fire products.
\newblock \emph{Remote Sensing of Environment}, 178:31--41, 2016.

\bibitem{hersbach2020}
H.~Hersbach et~al.
\newblock The ERA5 global reanalysis.
\newblock \emph{Quarterly Journal of the Royal Meteorological Society}, 146(730):1999--2049, 2020.

\bibitem{platt1999}
J.~Platt.
\newblock Probabilistic outputs for support vector machines and comparisons to regularized likelihood
methods.
\newblock In \emph{Advances in Large Margin Classifiers}, pp.~61--74. MIT Press, 1999.

\bibitem{schroeder2014}
W.~Schroeder, P.~Oliva, L.~Giglio, and I.~A. Csiszar.
\newblock The New VIIRS 375\,m active fire detection data product: algorithm description and initial
assessment.
\newblock \emph{Remote Sensing of Environment}, 143:85--96, 2014.

\bibitem{h3}
Uber Technologies.
\newblock H3: a hexagonal hierarchical geospatial indexing system.
\newblock \url{https://h3geo.org}, accessed November 2026.

\bibitem{vanwagnerpickett1985}
C.~E. Van~Wagner and T.~L. Pickett.
\newblock Equations and FORTRAN program for the Canadian Forest Fire Weather Index System.
\newblock Forestry Technical Report 33, Canadian Forestry Service, Ottawa, 1985.

\bibitem{vanwagner1987}
C.~E. Van~Wagner.
\newblock Development and structure of the Canadian Forest Fire Weather Index System.
\newblock Forestry Technical Report 35, Canadian Forestry Service, Ottawa, 1987.

\bibitem{wilks2011}
D.~S. Wilks.
\newblock \emph{Statistical Methods in the Atmospheric Sciences}, 3rd edition.
\newblock Academic Press, 2011.

\end{thebibliography}

\end{document}
